\documentclass[11pt,a4paper]{article}

\usepackage[utf8]{inputenc}
\usepackage[T1]{fontenc}
\usepackage{lmodern}
\usepackage{amsmath,amssymb,amsthm}
\usepackage{booktabs}
\usepackage{array}
\usepackage{caption}
\usepackage{geometry}
\usepackage{hyperref}
\usepackage{listings}
\usepackage{xcolor}
\usepackage{microtype}

\geometry{margin=1in}
\setlength{\emergencystretch}{3em}

\captionsetup[lstlisting]{singlelinecheck=false,justification=raggedright}

\hypersetup{
  colorlinks=true,
  linkcolor=blue!60!black,
  citecolor=green!50!black,
  urlcolor=blue!60!black
}

\lstdefinelanguage{Lean4}{
  morekeywords={def,theorem,lemma,abbrev,by,rfl,let,in,if,then,else,match,with,fun,true,false,set_option},
  sensitive=true,
  morecomment=[l]{--},
  morecomment=[s]{/-}{-/},
  morestring=[b]"
}

\lstset{
  language=Lean4,
  basicstyle=\ttfamily\footnotesize,
  keywordstyle=\color{blue!70!black}\bfseries,
  commentstyle=\color{gray!70!black}\itshape,
  stringstyle=\color{green!50!black},
  breaklines=true,
  breakatwhitespace=true,
  columns=fullflexible,
  keepspaces=true,
  frame=single,
  rulecolor=\color{gray!35},
  backgroundcolor=\color{gray!5},
  xleftmargin=0.5em,
  xrightmargin=0.5em,
  aboveskip=0.8em,
  belowskip=0.8em,
  literate={⟨}{{$\langle$}}1 {⟩}{{$\rangle$}}1 {×}{{$\times$}}1 {→}{{$\to$}}1 {∧}{{$\land$}}1 {·}{{$\cdot$}}1
}

\newtheorem{theorem}{Theorem}[section]
\newtheorem{proposition}[theorem]{Proposition}
\newtheorem{definition}[theorem]{Definition}
\newtheorem{remark}[theorem]{Remark}

\title{\textbf{Constructive Formalization of Shioda--Aoki Fermat Fourfold Character Cohomology in Lean~4}}

\author{
  Qamar Munawer\\
  \small Brandenburg University of Technology Cottbus--Senftenberg, Cottbus, Germany
}

\date{}

\begin{document}
\maketitle

\begin{abstract}
This paper presents a constructive, axiom-free formalization in Lean~4 of the Shioda--Aoki cyclotomic character classification governing primitive middle cohomology on Fermat fourfolds. The general Hodge Conjecture remains open, and this work formalizes strictly the decidable character combinatorics established under the classical theorems of Griffiths (1969), Shioda (1979), and Aoki (1983). Using fuel-bounded greatest common divisor evaluation and explicit six-tuple partition predicates with no external library dependencies, the Lean~4 kernel verifies that indecomposable residual characters vanish identically for all degrees $3 \le m \le 10$, decomposes the $1{,}751$ primitive Hodge classes of the sextic Fermat fourfold into $1{,}001 + 630 + 120 + 0$ alongside their exact $15$-matching multiplicity distribution, and isolates the four paired-diagonal $S_6$-orbits totaling $360$ residual characters at degree $m = 12$. All headline theorems compile via definitional kernel reduction and depend on zero axioms.
\end{abstract}

\section{Introduction and Epistemic Scope}

Let $X_m^n \subset \mathbb{P}^{n+1}$ denote the smooth complex projective Fermat hypersurface of dimension $n$ and degree $m \ge 3$, defined by the homogeneous equation
\begin{equation}
  x_0^m + x_1^m + \cdots + x_{n+1}^m = 0.
\end{equation}
When $n = 2p$ is even, the Hodge Conjecture for $X_m^n$ asks whether every rational cohomology class in $H^{2p}(X_m^n, \mathbb{Q}) \cap H^{p,p}(X_m^n)$ lies in the $\mathbb{Q}$-linear span of fundamental classes of algebraic subvarieties of codimension $p$~\cite{Deligne2000}. Although the general Hodge Conjecture for smooth complex projective varieties remains open, Fermat hypersurfaces admit a diagonal automorphism group action $(\mu_m)^{n+2}/\mu_m$ that decomposes the primitive middle cohomology $H^n_{\mathrm{prim}}(X_m^n, \mathbb{C})$ into one-dimensional eigenspaces indexed by characters $\alpha = (a_0, \dots, a_{n+1}) \in (\mathbb{Z}/m\mathbb{Z})^{n+2}$.

By Griffiths' residue theorem~\cite{Griffiths1969} and Shioda's character criterion~\cite{Shioda1979}, a character $\alpha$ with $1 \le a_i \le m - 1$ belongs to the rational Hodge subspace $H^{p,p}_{\mathrm{prim}}(X_m^n) \cap H^n(X_m^n, \mathbb{Q})$ if and only if every unit $t \in (\mathbb{Z}/m\mathbb{Z})^\times$ satisfies a constant fractional-part weight condition. Shioda~\cite{Shioda1979} proved that whenever $\alpha$ contains a complementary pair $a_i + a_j \equiv 0 \pmod{m}$, the corresponding Hodge class is algebraic via inductive join constructions from lower-dimensional Fermat varieties. Aoki~\cite{Aoki1983} extended this inductive spanning to three-element zero-sum triples and analyzed the remaining indecomposable characters.

Despite the foundational role of Fermat varieties in algebraic geometry, a prior-art scan across \texttt{mathlib4}, Google DeepMind's \texttt{formal-conjectures} repository, and the \texttt{FLT} formalization project confirms zero existing Lean~4 definitions or proofs of Shioda's Galois weight criterion, Aoki's four-band character partition, or Fermat fourfold Hodge character enumeration. The present formalization fills this gap constructively inside a standalone Lean~4.35.0 project (\texttt{ShiodaFormalConjectures.lean} and \texttt{HodgeVerification.lean}) with zero external library imports.

\begin{remark}[Epistemic Boundary]
  The formal theorems verified in Lean~4 establish exact finite combinatorial properties of the cyclotomic character set $\mathfrak{B}_m$ for fixed degrees $m$. They rely on the classical residue isomorphism~\cite{Griffiths1969} and inductive cycle theorems~\cite{Shioda1979,Aoki1983} as mathematical context and do not assert a proof of the general Hodge Conjecture.
\end{remark}

\section{Cyclotomic Character Group and the Griffiths--Shioda Galois Weight Criterion}

Fix $n = 4$ (fourfolds in $\mathbb{P}^5$) and degree $m \ge 3$. Let $\mathfrak{A}_m$ denote the primitive character set
\begin{equation}
  \mathfrak{A}_m = \left\{ \alpha = (a_0, a_1, a_2, a_3, a_4, a_5) \in \{1, \dots, m-1\}^6 \;\middle|\; \sum_{i=0}^5 a_i \equiv 0 \pmod{m} \right\}.
\end{equation}
For each unit $t \in (\mathbb{Z}/m\mathbb{Z})^\times$, the normalized Galois weight of $\alpha$ is defined by
\begin{equation}
  |t \cdot \alpha| = \frac{1}{m} \sum_{i=0}^5 \bigl((t a_i) \bmod m\bigr),
\end{equation}
where $(x \bmod m)$ denotes the canonical representative in $\{0, 1, \dots, m-1\}$. Since $1 \le a_i \le m-1$ and $\gcd(t, m) = 1$, each residue $(t a_i) \bmod m$ lies strictly in $\{1, \dots, m-1\}$, and the fractional part $\langle t a_i / m \rangle$ equals $((t a_i) \bmod m) / m$. Under Griffiths' residue identification between $H^{4-q, q}_{\mathrm{prim}}(X_m^4)$ and the Jacobian ring slice $R^{m(q+1)-6}(F_0)$, the middle $(2,2)$-Hodge characters $\mathfrak{B}_m \subset \mathfrak{A}_m$ are characterized by Shioda's theorem~\cite{Shioda1979}:
\begin{equation}
  \alpha \in \mathfrak{B}_m \iff \forall t \in (\mathbb{Z}/m\mathbb{Z})^\times, \quad \sum_{i=0}^5 \bigl((t a_i) \bmod m\bigr) = 3m.
\end{equation}

To guarantee that every theorem in Lean~4 evaluates by kernel reduction (\texttt{by rfl}) without invoking \texttt{propext} or well-founded recursion axioms from \texttt{Nat.gcd}, the coprime unit generator uses a fuel-bounded Euclidean algorithm \texttt{gcdFuel} and single-level structural pattern matching \texttt{with6Nat}. For any $1 \le t < m$, the Euclidean step $(a, b) \mapsto (b, a \bmod b)$ strictly decreases $b$ at every iteration, so $m$ fuel steps suffice to reach $b = 0$ and return $\gcd(t, m)$. Similarly, \texttt{with6Nat} decomposes a list of length six via nested one-step \texttt{match} expressions so that Lean's pattern compiler introduces no propositional extensionality equations. The verbatim Lean~4 definitions from \texttt{ShiodaFormalConjectures.lean} are reproduced below.

\begin{lstlisting}[caption={Constructive unit enumeration and Shioda $(2,2)$ Galois weight predicate.}]
def rangeList (start count : Nat) : List Nat :=
  match count with
  | 0 => []
  | c + 1 => start :: rangeList (start + 1) c

def gcdFuel : Nat -> Nat -> Nat -> Nat
  | 0, a, _ => a
  | k + 1, a, b => if b == 0 then a else gcdFuel k b (a % b)

/-- All units t in {1, ..., m - 1} coprime to m. -/
def coprimeUnits (m : Nat) : List Nat :=
  (rangeList 1 (m - 1)).filter (fun t => gcdFuel m t m == 1)

/-- Shioda (1979) / Griffiths residue criterion for a character alpha = [a0, ..., a5]
    to lie in H^{2,2}_{prim}(X_m^4) \cap H^4(X_m^4, Q):
    a.length == 6, every 1 <= a_i <= m - 1, and for EVERY t in (Z/mZ)^x,
    sum_{i=0}^5 ((t * a_i) mod m) = 3 * m. -/
def isShiodaHodge22 (m : Nat) (a : List Nat) : Bool :=
  a.length == 6 &&
  a.all (fun x => Nat.blt 0 x && Nat.blt x m) &&
  (coprimeUnits m).all (fun t =>
    (a.foldl (fun acc x => acc + (t * x) % m) 0) == 3 * m)

def with6Nat (a : List Nat) (f : Nat -> Nat -> Nat -> Nat -> Nat -> Nat -> Bool) : Bool :=
  match a with
  | [] => false
  | a0 :: r0 => match r0 with
    | [] => false
    | a1 :: r1 => match r1 with
      | [] => false
      | a2 :: r2 => match r2 with
        | [] => false
        | a3 :: r3 => match r3 with
          | [] => false
          | a4 :: r4 => match r4 with
            | [] => false
            | a5 :: r5 => match r5 with
              | [] => f a0 a1 a2 a3 a4 a5
              | _ :: _ => false
\end{lstlisting}

Observe that checking every unit $t \in (\mathbb{Z}/m\mathbb{Z})^\times$ is essential. For instance, at $m = 12$, the tuple $(1, 1, 1, 5, 7, 9)$ sums to $24 = 2m$ at $t = 1$ (and $48 = 4m$ at $t = 7$), placing it in $H^{3,1} \oplus H^{1,3}$ rather than $H^{2,2}$, even though its average weight across units is $36$. Requiring equality to $3m$ universally across \texttt{coprimeUnits m} eliminates all such spurious candidates.

\section{Algorithmic Four-Band Partition and Multinomial \texorpdfstring{$S_6$}{S6}-Orbit Compression}

Following Shioda~\cite{Shioda1979} and Aoki~\cite{Aoki1983}, any Hodge character $\alpha \in \mathfrak{B}_m$ is classified into one of four mutually disjoint geometric bands according to zero-sum partitions modulo $m$ among its six indices $\{0, 1, 2, 3, 4, 5\}$:

\begin{enumerate}
  \item \textbf{Band 1 (\texttt{isBand1CompletePairs}):} The six indices partition into three disjoint pairs $\{i_1, j_1\}, \{i_2, j_2\}, \{i_3, j_3\}$ satisfying $a_{i_k} + a_{j_k} = m$ for $k = 1, 2, 3$. Geometrically, these characters correspond to classes supported on linear $2$-planes inside $X_m^4$.
  \item \textbf{Band 2 (\texttt{isBand2ThreePlusThree}):} $\alpha$ is not in Band~1, and the six indices partition into two disjoint triples whose coordinate sums are $\{m, 2m\}$. These correspond to $3+3$ inductive Fermat curve joins.
  \item \textbf{Band 3 (\texttt{isBand3TwoPlusFour}):} $\alpha$ is neither in Band~1 nor in Band~2, and there exists at least one pair $i < j$ with $a_i + a_j = m$. The complementary four coordinates sum to $2m$, corresponding to $2+4$ surface-curve joins via Lefschetz on the Fermat surface $X_m^2 \subset \mathbb{P}^3$.
  \item \textbf{Band 4 (\texttt{isBand4Residual}):} Indecomposable characters in $\mathfrak{B}_m$ belonging to none of Bands~1, 2, or~3.
\end{enumerate}

By construction, the conditional guards in the second and third band predicates enforce mutual disjointness, so every $\alpha \in \mathfrak{B}_m$ satisfies exactly one of the four predicates. Using the six-coordinate unpacking helper \texttt{with6Nat}, the Lean~4 definitions check all $15$ pair-matchings, $10$ triple bipartitions, and $15$ index pairs.

\begin{lstlisting}[caption={Four-band geometric partition predicates in \texttt{ShiodaFormalConjectures.lean}.}]
/-- Band 1: True iff the 6 indices {0..5} partition into 3 disjoint pairs summing to m. -/
def isBand1CompletePairs (m : Nat) (a : List Nat) : Bool :=
  with6Nat a fun a0 a1 a2 a3 a4 a5 =>
    (a0 + a1 == m && a2 + a3 == m && a4 + a5 == m) ||
    (a0 + a1 == m && a2 + a4 == m && a3 + a5 == m) ||
    (a0 + a1 == m && a2 + a5 == m && a3 + a4 == m) ||
    (a0 + a2 == m && a1 + a3 == m && a4 + a5 == m) ||
    (a0 + a2 == m && a1 + a4 == m && a3 + a5 == m) ||
    (a0 + a2 == m && a1 + a5 == m && a3 + a4 == m) ||
    (a0 + a3 == m && a1 + a2 == m && a4 + a5 == m) ||
    (a0 + a3 == m && a1 + a4 == m && a2 + a5 == m) ||
    (a0 + a3 == m && a1 + a5 == m && a2 + a4 == m) ||
    (a0 + a4 == m && a1 + a2 == m && a3 + a5 == m) ||
    (a0 + a4 == m && a1 + a3 == m && a2 + a5 == m) ||
    (a0 + a4 == m && a1 + a5 == m && a2 + a3 == m) ||
    (a0 + a5 == m && a1 + a2 == m && a3 + a4 == m) ||
    (a0 + a5 == m && a1 + a3 == m && a2 + a4 == m) ||
    (a0 + a5 == m && a1 + a4 == m && a2 + a3 == m)

/-- Band 2: True iff not Band 1 AND the 6 indices partition into two disjoint triples
    whose sums are {m, 2*m}. -/
def isBand2ThreePlusThree (m : Nat) (a : List Nat) : Bool :=
  if isBand1CompletePairs m a then false
  else
    with6Nat a fun a0 a1 a2 a3 a4 a5 =>
      let check (s1 s2 : Nat) : Bool := (s1 == m && s2 == 2 * m) || (s1 == 2 * m && s2 == m)
      check (a0 + a1 + a2) (a3 + a4 + a5) ||
      check (a0 + a1 + a3) (a2 + a4 + a5) ||
      check (a0 + a1 + a4) (a2 + a3 + a5) ||
      check (a0 + a1 + a5) (a2 + a3 + a4) ||
      check (a0 + a2 + a3) (a1 + a4 + a5) ||
      check (a0 + a2 + a4) (a1 + a3 + a5) ||
      check (a0 + a2 + a5) (a1 + a3 + a4) ||
      check (a0 + a3 + a4) (a1 + a2 + a5) ||
      check (a0 + a3 + a5) (a1 + a2 + a4) ||
      check (a0 + a4 + a5) (a1 + a2 + a3)

/-- Band 3: True iff not Band 1, not Band 2, AND there exists at least one pair i < j with a_i + a_j = m. -/
def isBand3TwoPlusFour (m : Nat) (a : List Nat) : Bool :=
  if isBand1CompletePairs m a || isBand2ThreePlusThree m a then false
  else
    with6Nat a fun a0 a1 a2 a3 a4 a5 =>
      (a0 + a1 == m) || (a0 + a2 == m) || (a0 + a3 == m) || (a0 + a4 == m) || (a0 + a5 == m) ||
      (a1 + a2 == m) || (a1 + a3 == m) || (a1 + a4 == m) || (a1 + a5 == m) ||
      (a2 + a3 == m) || (a2 + a4 == m) || (a2 + a5 == m) ||
      (a3 + a4 == m) || (a3 + a5 == m) ||
      (a4 + a5 == m)

/-- Band 4: Residual Hodge (2,2) characters not in Bands 1, 2, or 3. -/
def isBand4Residual (m : Nat) (a : List Nat) : Bool :=
  isShiodaHodge22 m a &&
  !(isBand1CompletePairs m a) &&
  !(isBand2ThreePlusThree m a) &&
  !(isBand3TwoPlusFour m a)
\end{lstlisting}

Direct enumeration of $(m-1)^6$ ordered tuples becomes computationally expensive inside a proof assistant kernel as $m$ grows (exceeding $1.77 \times 10^6$ tuples at $m = 12$). Because both \texttt{isShiodaHodge22} and all four band predicates are invariant under the symmetric group $S_6$, the search space compresses identically to non-decreasing sorted representatives $1 \le a_0 \le a_1 \le a_2 \le a_3 \le a_4 \le a_5 \le m-1$ satisfying $\sum a_i = 3m$, weighted by the exact multinomial orbit cardinality
\begin{equation}
  |S_6 \cdot \alpha| = \frac{6!}{\prod_{x} (\operatorname{mult}(x, \alpha))!}.
\end{equation}
For example, a tuple with six distinct entries has orbit size $720$, a paired-diagonal tuple $(a, a, b, b, c, c)$ with $a < b < c$ has orbit size $720 / (2!2!2!) = 90$, and the constant tuple $(m/2, \dots, m/2)$ has orbit size $720 / 6! = 1$. The sorted generator \texttt{sortedHodge22Tuples} prunes impossible branches early via $a_0 \le m/2$ and partial-sum monotonicity bounds $a_0 + \dots + a_{k-1} + (6-k)a_k \le 3m$, and \texttt{summarizeHodgeTuples} accumulates all band counts and Band~4 representatives in a single fold.

\begin{lstlisting}[caption={Multinomial $S_6$-orbit size, sorted enumeration, and single-pass census.}]
def fact : Nat -> Nat
  | 0 => 1
  | n + 1 => (n + 1) * fact n

def countOcc (x : Nat) : List Nat -> Nat
  | [] => 0
  | y :: ys => (if x == y then 1 else 0) + countOcc x ys

def uniqueList : List Nat -> List Nat
  | [] => []
  | x :: xs =>
    let u := uniqueList xs
    if u.contains x then u else x :: u

/-- Exact S_6-orbit size: 6! / prod(c_k!) for a 6-tuple a. -/
def s6OrbitSize (a : List Nat) : Nat :=
  let denom := (uniqueList a).foldl (fun acc x => acc * fact (countOcc x a)) 1
  fact a.length / denom

def sortedHodge22Tuples (m : Nat) : List (List Nat) :=
  let units := coprimeUnits m
  let checkHodge (a0 a1 a2 a3 a4 a5 : Nat) : Bool :=
    units.all (fun t =>
      ((t * a0) % m + (t * a1) % m + (t * a2) % m +
       (t * a3) % m + (t * a4) % m + (t * a5) % m) == 3 * m)
  let vals0 := rangeList 1 (m / 2)
  vals0.foldl (fun acc0 a0 =>
    let vals1 := rangeList a0 (m - a0)
    acc0 ++ vals1.foldl (fun acc1 a1 =>
      if Nat.blt (3 * m) (a0 + 5 * a1) then acc1 else
      let vals2 := rangeList a1 (m - a1)
      acc1 ++ vals2.foldl (fun acc2 a2 =>
        if Nat.blt (3 * m) (a0 + a1 + 4 * a2) then acc2 else
        let vals3 := rangeList a2 (m - a2)
        acc2 ++ vals3.foldl (fun acc3 a3 =>
          if Nat.blt (3 * m) (a0 + a1 + a2 + 3 * a3) then acc3 else
          let vals4 := rangeList a3 (m - a3)
          acc3 ++ vals4.foldl (fun acc4 a4 =>
            let s := a0 + a1 + a2 + a3 + a4
            if Nat.blt (3 * m) (s + a4) then acc4 else
            let a5 := 3 * m - s
            if Nat.blt a5 m && checkHodge a0 a1 a2 a3 a4 a5 then
              acc4 ++ [[a0, a1, a2, a3, a4, a5]]
            else acc4
          ) []
        ) []
      ) []
    ) []
  ) []

def summarizeHodgeTuples (m : Nat) (ts : List (List Nat)) :
    Nat × Nat × Nat × Nat × Nat × List (List Nat) :=
  ts.foldl (fun (h22, b1, b2, b3, b4, reps) tuple =>
    let w := s6OrbitSize tuple
    if isBand1CompletePairs m tuple then (h22 + w, b1 + w, b2, b3, b4, reps)
    else if isBand2ThreePlusThree m tuple then (h22 + w, b1, b2 + w, b3, b4, reps)
    else if isBand3TwoPlusFour m tuple then (h22 + w, b1, b2, b3 + w, b4, reps)
    else (h22 + w, b1, b2, b3, b4 + w, reps ++ [tuple])
  ) (0, 0, 0, 0, 0, [])

def fermatCensusSummary (m : Nat) : Nat × Nat × Nat × Nat × Nat × List (List Nat) :=
  summarizeHodgeTuples m (sortedHodge22Tuples m)
\end{lstlisting}

\section{The Sextic Calabi--Yau Fourfold (\texorpdfstring{$m = 6$}{m = 6}) and the 15-Matching Incidence Distribution}

For degree $m = 6$, $X_6^4 \subset \mathbb{P}^5$ is a Calabi--Yau fourfold with Milnor number $(6-1)^6 = 15{,}625$ and primitive Hodge dimension $h^{2,2}_{\mathrm{prim}}(X_6^4) = 1{,}751$. Evaluating \texttt{fermatCensusSummary 6} in Lean~4 proves that the four-band decomposition sums exactly to $1{,}751$, and that all Band~3 characters universally satisfy the shifted $2+4$ Lefschetz split in $R^8(F_{\mathbb{P}^3})$ (where shifting each coordinate $a_i \mapsto a_i - 1 \in \{0, \dots, 4\}$ maps $\sum a_i = 18$ to $\sum (a_i - 1) = 12$ and $a_i + a_j = 6$ to $(a_i - 1) + (a_j - 1) = 4$).

\begin{lstlisting}[caption={Sextic Calabi--Yau fourfold census and Band~3 universal split.}]
theorem shioda_m6_orbit_census_cert :
    (5 : Nat) ^ 6 = 15625 /\
    fermatCensusSummary 6 = (1751, 1001, 630, 120, 0, []) :=
  ⟨by rfl, by rfl⟩

def fermat_four_band_total : Nat :=
  fermat_band1_count 6 + fermat_band2_count 6 + fermat_band3_count 6 + fermat_band4_count 6

theorem fermat_four_band_closure_eq : fermat_four_band_total = h22_prim := by rfl

def has_pair_sum_4 : List Nat -> Bool
  | [] => false
  | x :: xs => (xs.any (fun y => x + y == 4)) || has_pair_sum_4 xs

def has_2_plus_4_split (a : List Nat) : Bool :=
  (a.length == 6) && (a.foldl (. + .) 0 == 12) && (a.all (fun x => x <= 4)) && has_pair_sum_4 a

def band3_sorted_tuples_m6 : List (List Nat) :=
  (sortedHodge22Tuples 6).filter (isBand3TwoPlusFour 6)

theorem band3_partitions_satisfy_2_plus_4 :
  (band3_sorted_tuples_m6.length > 0 &&
   band3_sorted_tuples_m6.all (fun a => has_2_plus_4_split (a.map (fun x => x - 1)))) = true := by rfl
\end{lstlisting}

On the sextic Fermat fourfold, each of the $15$ perfect matchings of $\{0, 1, 2, 3, 4, 5\}$ generates a coordinate $2$-plane block supporting $5^3 = 125$ characters, for a total block-incidence sum of $15 \times 125 = 1{,}875$ across the $1{,}001$ distinct Band~1 characters. By generating all $15$ matchings algorithmically in \texttt{all\_perfect\_matchings\_6} and partitioning \texttt{band1\_sorted\_tuples\_m6} by the number $k$ of matchings satisfied by each character, Lean~4 verifies the exact multiplicity distribution
\begin{equation}
  360 \cdot 1 + 540 \cdot 2 + 60 \cdot 3 + 40 \cdot 6 + 1 \cdot 15 = 1{,}875, \qquad 360 + 540 + 60 + 40 + 1 = 1{,}001.
\end{equation}
Geometrically, the unique character of multiplicity $k = 15$ is $(3, 3, 3, 3, 3, 3)$, which satisfies every perfect matching because $3 + 3 = 6$ for all pairs. The $40$ characters of multiplicity $k = 6$ arise from permutations of $(1, 1, 1, 5, 5, 5)$ and $(2, 2, 2, 4, 4, 4)$, which admit $3! = 6$ bijections between their two complementary triples.

\begin{table}[htbp]
  \centering
  \caption{Exact multiplicity partition of the $1{,}001$ Band~1 characters across the $15$ perfect matchings of $\{0,1,2,3,4,5\}$ at degree $m = 6$.}
  \label{tab:band1_mult}
  \begin{tabular}{cccc}
    \toprule
    Multiplicity $k$ & Character Count $n_k$ & Block Incidences $k \cdot n_k$ & Representative Orbit Shape \\
    \midrule
    $1$ & $360$ & $360$ & $(1, 5, 2, 4, 3, 3)$ \\
    $2$ & $540$ & $1{,}080$ & $(1, 1, 5, 5, 2, 4)$ \\
    $3$ & $60$ & $180$ & $(1, 2, 3, 3, 4, 5)$ \\
    $6$ & $40$ & $240$ & $(1, 1, 1, 5, 5, 5)$ \\
    $15$ & $1$ & $15$ & $(3, 3, 3, 3, 3, 3)$ \\
    \midrule
    Total & $1{,}001$ & $1{,}875$ & $15 \times 125$ block incidences \\
    \bottomrule
  \end{tabular}
\end{table}

\begin{lstlisting}[caption={Algorithmic $15$-matching multiplicity theorem in \texttt{HodgeVerification.lean}.}]
def all_perfect_matchings_6 : List (List (Nat × Nat)) :=
  [1, 2, 3, 4, 5].foldl (fun acc b =>
    match [1, 2, 3, 4, 5].filter (fun x => x != b) with
    | [] => acc
    | r0 :: t0 => match t0 with
      | [] => acc
      | r1 :: t1 => match t1 with
        | [] => acc
        | r2 :: t2 => match t2 with
          | [] => acc
          | r3 :: t3 => match t3 with
            | [] =>
              acc ++ [
                [(0, b), (r0, r1), (r2, r3)],
                [(0, b), (r0, r2), (r1, r3)],
                [(0, b), (r0, r3), (r1, r2)]
              ]
            | _ :: _ => acc
  ) []

def band1_sorted_tuples_m6 : List (List Nat) :=
  (sortedHodge22Tuples 6).filter (isBand1CompletePairs 6)

def band1_count_by_multiplicity (k : Nat) : Nat :=
  (band1_sorted_tuples_m6.filter (fun a => matchingSatisfactionCount 6 a == k)).foldl
    (fun acc a => acc + s6OrbitSize a) 0

def band1_multiplicity_pairs : List (Nat × Nat) :=
  [1, 2, 3, 6, 15].map (fun k => (k, band1_count_by_multiplicity k))

set_option maxRecDepth 4096 in
theorem fifteen_matching_band1_block_incidence_eq :
  (fifteen_perfect_matchings == 15 &&
   modes_per_matching_block == 125 &&
   band1_block_total_incidences == 1875 &&
   band1_mult_mode_sum == fermat_band1_count 6 &&
   band1_mult_incidence_sum == band1_block_total_incidences &&
   band1_block_total_incidences - fermat_band1_count 6 == 874 &&
   band1_multiplicity_pairs == [(1, 360), (2, 540), (3, 60), (6, 40), (15, 1)]) = true := by rfl
\end{lstlisting}

\section{Band~4 Vanishing for \texorpdfstring{$3 \le m \le 10$}{3 <= m <= 10} and the Degree \texorpdfstring{$m = 12$}{m = 12} Galois Orbit}

Evaluating \texttt{fermat\_band4\_count} across all degrees $m \in \{3, 4, 5, 6, 7, 8, 9, 10\}$ establishes constructively that Band~4 is identically empty throughout this range, confirming Aoki's low-degree vanishing theorem~\cite{Aoki1983} for fourfolds. At $m = 12$, the unit group is $(\mathbb{Z}/12\mathbb{Z})^\times = \{1, 5, 7, 11\}$. The kernel-verified census finds $52{,}991$ primitive Hodge $(2,2)$ characters decomposed into $15{,}101$ (Band~1), $21{,}690$ (Band~2), $15{,}840$ (Band~3), and $360$ residual Band~4 characters.

\begin{lstlisting}[caption={Low-degree Band~4 vanishing and $m = 12$ Galois orbit census.}]
theorem shioda_aoki_fermat_fourfold_census_m3_to_m10_empty_band4 :
    [3, 4, 5, 6, 7, 8, 9, 10].map fermat_band4_count = [0, 0, 0, 0, 0, 0, 0, 0] := by rfl

set_option maxRecDepth 4096 in
theorem shioda_aoki_fermat_m12_census_cert :
    fermatCensusSummary 12 =
      (52991, 15101, 21690, 15840, 360,
        [[1, 1, 8, 8, 9, 9], [3, 3, 4, 4, 11, 11], [3, 3, 7, 7, 8, 8], [4, 4, 5, 5, 9, 9]]) :=
  by rfl
\end{lstlisting}

Remarkably, the four sorted Band~4 representatives at $m = 12$ form a single closed $(\mathbb{Z}/12\mathbb{Z})^\times$-orbit. Starting from $\alpha_1 = (1, 1, 8, 8, 9, 9)$, multiplying modulo $12$ by the remaining units and sorting yields:
\begin{align}
  5 \cdot (1, 1, 8, 8, 9, 9) \bmod 12 &= (5, 5, 4, 4, 9, 9) \sim (4, 4, 5, 5, 9, 9), \\
  7 \cdot (1, 1, 8, 8, 9, 9) \bmod 12 &= (7, 7, 8, 8, 3, 3) \sim (3, 3, 7, 7, 8, 8), \\
  11 \cdot (1, 1, 8, 8, 9, 9) \bmod 12 &= (11, 11, 4, 4, 3, 3) \sim (3, 3, 4, 4, 11, 11).
\end{align}
Each representative has three distinct entries of multiplicity $2$, giving $S_6$-orbit size $6! / (2!2!2!) = 90$ and total Band~4 cardinality $4 \times 90 = 360$. Because every entry occurs in an equal pair $(a, a, b, b, c, c)$, all $360$ classes are algebraic via paired-diagonal pullback~\cite{Aoki1983}.

\begin{table}[htbp]
  \centering
  \caption{Lean~4 kernel-verified primitive $(2,2)$-Hodge character census on $X_m^4 \subset \mathbb{P}^5$ for $3 \le m \le 10$ and $m = 12$.}
  \label{tab:shioda_census}
  \begin{tabular}{ccccccc}
    \toprule
    Degree $m$ & Total $\mathfrak{B}_m$ & Band 1 & Band 2 & Band 3 & Band 4 & Band 4 Representatives \\
    \midrule
    $3 \le m \le 10$ & --- & --- & --- & --- & $0$ & $\emptyset$ \\
    $6$ & $1{,}751$ & $1{,}001$ & $630$ & $120$ & $0$ & $\emptyset$ \\
    $12$ & $52{,}991$ & $15{,}101$ & $21{,}690$ & $15{,}840$ & $360$ & $4$ paired orbits ($\times 90$) \\
    \bottomrule
  \end{tabular}
\end{table}

\section{Constructive Complexity and Symmetric-Group Orbit Enumeration}

A central requirement of kernel-level certification in Lean~4 is that every decision procedure must terminate structurally without invoking well-founded recursion lemmas or classical choice axioms, and must execute within the kernel's reduction budget. Two design choices make the degree $m \le 12$ census tractable by pure definitional reduction (\texttt{by rfl}).

\subsection{Fuel-Bounded Euclidean Reduction}
Standard Euclidean division in Lean's core library (\texttt{Nat.gcd}) uses well-founded recursion on the second argument, which introduces a dependency on \texttt{propext} in the compiled proof term. Replacing \texttt{Nat.gcd} with \texttt{gcdFuel m t m} bounds the recursion depth structurally by $m$. Because each Euclidean step satisfies $a \bmod b < b \le m$, the second argument reaches $0$ in at most $m$ steps for all $t \in \{1, \dots, m-1\}$. Consequently, \texttt{coprimeUnits m} is entirely free of axioms while remaining extensionally equal to $\{t \in \{1, \dots, m-1\} \mid \gcd(t, m) = 1\}$.

\subsection{Monotone Generation and Multinomial Orbit Weights}
Evaluating \texttt{isShiodaHodge22} naively over all $(m-1)^6$ ordered tuples at $m = 12$ would require testing $11^6 = 1{,}771{,}561$ candidates inside the Lean kernel. Instead, \texttt{sortedHodge22Tuples} generates only non-decreasing tuples $1 \le a_0 \le a_1 \le a_2 \le a_3 \le a_4 \le a_5 \le m - 1$, reducing the candidate space at $m = 12$ to $\binom{16}{6} = 8{,}008$ tuples---a $221\times$ speedup. Moreover, the sixth coordinate $a_5$ is uniquely determined at $t = 1$ by $a_5 = 3m - (a_0 + a_1 + a_2 + a_3 + a_4)$, which directly determines $a_5$ in the innermost loop without a sixth nested enumeration.

Each non-decreasing tuple $\alpha$ represents an entire $S_6$-orbit of ordered characters, all of which share the same Galois-orbit sum property and band membership. The cardinality of this $S_6$-orbit depends solely on the multiplicity partition $\lambda \vdash 6$ of repeated coordinate values in $\alpha$, as computed by \texttt{s6OrbitSize} via $6! / \prod_v c_v!$. Table~\ref{tab:s6_partitions} lists all $11$ integer partitions of $6$ together with their stabilizer orders and orbit weights.

\begin{table}[htbp]
  \centering
  \caption{All $11$ coordinate multiplicity partitions $\lambda \vdash 6$ and their $S_6$-orbit weights computed by \texttt{s6OrbitSize}.}
  \label{tab:s6_partitions}
  \begin{tabular}{cccc}
    \toprule
    Multiplicity Partition $\lambda \vdash 6$ & Example Pattern & Stabilizer $\prod c_v!$ & $S_6$-Orbit Size \\
    \midrule
    $(1,1,1,1,1,1)$ & $(a,b,c,d,e,f)$ & $1$ & $720$ \\
    $(2,1,1,1,1)$ & $(a,a,b,c,d,e)$ & $2$ & $360$ \\
    $(2,2,1,1)$ & $(a,a,b,b,c,d)$ & $4$ & $180$ \\
    $(3,1,1,1)$ & $(a,a,a,b,c,d)$ & $6$ & $120$ \\
    $(2,2,2)$ & $(a,a,b,b,c,c)$ & $8$ & $90$ \\
    $(3,2,1)$ & $(a,a,a,b,b,c)$ & $12$ & $60$ \\
    $(4,1,1)$ & $(a,a,a,a,b,c)$ & $24$ & $30$ \\
    $(3,3)$ & $(a,a,a,b,b,b)$ & $36$ & $20$ \\
    $(4,2)$ & $(a,a,a,a,b,b)$ & $48$ & $15$ \\
    $(5,1)$ & $(a,a,a,a,a,b)$ & $120$ & $6$ \\
    $(6)$ & $(a,a,a,a,a,a)$ & $720$ & $1$ \\
    \bottomrule
  \end{tabular}
\end{table}

\appendix
\section{Verbatim Kernel Build and Zero-Axiom Verification Receipts}

All Lean~4 theorems presented in this paper were verified in a clean build using Lean~4.35.0 and Lake~5.0.0 via \texttt{lake clean \&\& lake build}. Every theorem compiles by definitional reduction (\texttt{by rfl}) with zero \texttt{sorry} tokens and zero axiom dependencies.

\begin{lstlisting}[language={},basicstyle=\ttfamily\scriptsize,caption={Raw \texttt{lake clean \&\& lake build} and \texttt{\#print axioms} output.}]
[2/5] Built ShiodaFormalConjectures (64s)
info: ShiodaFormalConjectures.lean:466:0: 'shioda_m6_orbit_census_cert' does not depend on any axioms
info: ShiodaFormalConjectures.lean:467:0: 'shioda_aoki_fermat_fourfold_census_m3_to_m10_empty_band4' does not depend on any axioms
info: ShiodaFormalConjectures.lean:468:0: 'shioda_aoki_fermat_m12_census_cert' does not depend on any axioms
info: ShiodaFormalConjectures.lean:469:0: 'weyl_d4_transition_matrix_invertible_cert' does not depend on any axioms
info: ShiodaFormalConjectures.lean:470:0: 'macaulay_plane_and_quadric_rank_witness_cert' does not depend on any axioms
info: ShiodaFormalConjectures.lean:471:0: 'macaulay_pairing_matrix_rank_cert' does not depend on any axioms
[3/5] Built HodgeVerification (22s)
info: HodgeVerification.lean:321:0: 'band3_partitions_satisfy_2_plus_4' does not depend on any axioms
info: HodgeVerification.lean:322:0: 'fermat_four_band_closure_eq' does not depend on any axioms
info: HodgeVerification.lean:701:0: 'fifteen_matching_band1_block_incidence_eq' does not depend on any axioms
[4/5] Built LeanDemoProject (523ms)
Build completed successfully (5 jobs).
\end{lstlisting}

\begin{thebibliography}{9}
\bibitem{Aoki1983}
N.~Aoki, \emph{On the Hodge conjecture for Fermat varieties}, Comment. Math. Univ. St. Pauli \textbf{32} (1983), 181--187.

\bibitem{Deligne2000}
P.~Deligne, \emph{The Hodge Conjecture}, Clay Mathematics Institute Millennium Prize Problems, 2000.

\bibitem{Griffiths1969}
P.~A.~Griffiths, \emph{On the periods of certain rational integrals} (Parts~1 and~2), Ann. of Math. (2) \textbf{90} (1969), 460--541.

\bibitem{Shioda1979}
T.~Shioda, \emph{The Hodge conjecture for Fermat varieties}, Math. Ann. \textbf{245} (1979), 175--184.
\end{thebibliography}

\end{document}
